% =====================================================================================================
% S-PAPER: THEORY SECTION SKELETON for the SMC-BVS paper (JRSS-B target).
% Self-contained article fragment: compile as is with pdflatex (amsmath, amssymb, amsthm only).
%
% DRAFT. Nothing here is reviewed or adopted. Every statement carries a comment line of the form
%     % STATUS: <author status>; VERIFIED-BY: <branches>; SOURCE: <branch:file:label>
% Status words are the campaign protocol's (PROVED-LOCKED, PROVED, PROVED-UNDER, PROOF SKETCH, CONJECTURE,
% OPEN, IMPL, IDEAL, COUNTEREXAMPLE) and are never upgraded here.  "VERIFIED-BY: W..." is a verifier of the SAME
% model family; "A-REV1/A-REV2" are cross-family (GPT-6 Astra) reviews, whose verdicts are quoted, never extended.
% "VERIFIED-BY: none" means no verifier run exists on the branches fetched when this was written.
% Statements are copied from their sources; where a source is a candidate proof, the statement is a candidate.
% No proofs: each result points to a Supplement section (placeholders "Supplement S.x") and to its source branch.
% Not incorporated: W6-FIXREV1 (A-REV1 repairs, new T38 files) and W6-FIXREV2 (A-REV2 repairs, STATEMENTS.md). Both branches had
% been pushed when this was finalised (7e58127, d213e7b) but contained only scripts/checks/notes: no T38_STATEMENT.md, no STATEMENTS.md,
% no REPORT.md. The statements below are therefore the pre-W6 versions plus the A-REV1/A-REV2 verdicts. See README.md, 'Pending edits'.
% =====================================================================================================
\documentclass[11pt]{article}
\usepackage[T1]{fontenc}
\usepackage{amsmath,amssymb,amsthm}
\usepackage[margin=1in]{geometry}

% ---------- theorem environments ----------
\newtheorem{theorem}{Theorem}
\newtheorem{thmpart}{Theorem}[theorem]
\renewcommand{\thethmpart}{\thetheorem(\roman{thmpart})}
\newtheorem{proposition}{Proposition}
\newtheorem{lemma}{Lemma}
\theoremstyle{definition}
\newtheorem{definition}{Definition}
\newtheorem{assumption}{Assumption}
\theoremstyle{remark}
\newtheorem{remark}{Remark}

% ---------- macros ----------
\renewcommand{\Pr}{\mathrm{P}}
\newcommand{\E}{\mathrm{E}}
\newcommand{\Sz}{S_0}
\newcommand{\Sh}{\widehat S}
\newcommand{\Spl}{\widehat S^{+}}
\newcommand{\Zh}{\widehat Z}
\newcommand{\PIPh}{\widehat{\mathrm{PIP}}}
\newcommand{\PIP}{\mathrm{PIP}}
\newcommand{\MPM}{\mathrm{MPM}}
\newcommand{\HPM}{\mathrm{HPM}}
\newcommand{\cC}{\mathcal C}
\newcommand{\cA}{\mathcal A}
\newcommand{\cS}{\mathcal S}
\newcommand{\cO}{\mathcal O}
\newcommand{\cE}{\mathcal E}
\newcommand{\cG}{\mathcal G}
\newcommand{\cR}{\mathcal R}
\newcommand{\Phidist}{\Phi}
\newcommand{\Lat}{\mathrm{Lat}}
\newcommand{\given}{\mid}
\newcommand{\Ind}{\mathbf 1}
\DeclareMathOperator*{\argmax}{arg\,max}
\DeclareMathOperator{\Bin}{Bin}

\title{Theory section (skeleton): consistency of the reported model, evidence and inclusion probabilities
for SMC-BVS}
\author{Draft for the JRSS-B paper --- S-PAPER, branch \texttt{theory/S-PAPER}}
\date{2 November 2026 --- DRAFT, unreviewed}

\begin{document}
\maketitle

\noindent\textbf{Reading guide.} Every numbered result below is a candidate statement assembled from the
verification campaign. The comment lines in the \LaTeX{} source give, for every statement, its author status, the
branches that verified it and the exact source label. In this typeset text, the phrase ``Proof: see Supplement
\S\,X (source branch \dots)'' stands for a proof that is \emph{not} reproduced here. Where the source is a
single-run proof not checked by a second run, this is said in a remark, and the status words of the campaign are
never upgraded.

% =====================================================================================================
\section{Setting and notation}\label{sec:setting}
% =====================================================================================================

\subsection{Data, target and decision objects}

% STATUS: definitions; VERIFIED-BY: n/a; SOURCE: main:Docs/Theory/theorems/00_COMMON.md:§1
Let $(X^r,y^r)$ be raw data, $n\ge 3$, with $X^r$ having no constant column and $y^r$ non-constant. With
$C=I-\mathbf 1\mathbf 1^\top/n$ and $D_x$ the diagonal matrix of retained raw column standard deviations, the
standardised data are $X=CX^rD_x^{-1}$ and $y=Cy^r/s_y$ with $s_y^2=\|Cy^r\|^2/(n-1)$, so that
$\|x_j\|^2=\|y\|^2=n-1$. The state space is $\Omega=\{S\subseteq[p]:|S|\le K\}$ with $K=\min(p,\lfloor n/2\rfloor)$.
For $S\in\Omega$ write $G_S=X_S^\top X_S=(n-1)R_S$, $z_S=X_S^\top y$, $H_S=(n+1)G_S+I$, $D_S=\det H_S/\det(G_S+I)$,
$Q_S=y^\top y-n\,z_S^\top H_S^{-1}z_S$ and $\nu=0.01+n/2$. The target is the Bayesian variable-selection posterior with
$\beta_S\sim N(0,n\sigma^2(G_S+I)^{-1})$, $\sigma^2\sim\mathrm{IG}(0.01,0.01)$ and a beta--binomial size prior
BB$(1,p)$ capped at $K$: the marginal likelihood is
$m(S)=(2b_0)^{a_0}\Gamma(\nu)\pi^{-n/2}\Gamma(a_0)^{-1}D_S^{-1/2}(2b_0+Q_S)^{-\nu}$ ($a_0=b_0=0.01$), the prior is
$\varpi(S)=w(|S|)/W_K$ with $w(k)=\tfrac12\binom{2p-1}{k}^{-1}$ and $W_K=\sum_{k\le K}\binom pk w(k)$, and
\[
F(S)=\log[\varpi(S)m(S)],\qquad Z=\sum_{S\in\Omega}\varpi(S)m(S),\qquad \pi(S)=e^{F(S)}/Z .
\]
Write $\PIP_j=\pi(j\in S)$ and $\MPM=\{j:\PIP_j>\tfrac12\}$. The \emph{global} HPM is $\argmax_\Omega F$ with ties
broken by smaller size, then lexicographically. The \emph{reported} HPM $\Sh$ is the same rule applied to the set
$\cS$ of all models ever scored.

\subsection{Truth and asymptotic regime}

% STATUS: definitions; VERIFIED-BY: n/a; SOURCE: main:Docs/Theory/theorems/00_COMMON.md:§1,§2 (R1), §3 (G, R)
The statistical model is $y^r=\alpha^*\mathbf 1+X^r\beta^*+\varepsilon$ with $\mu=CX^r\beta^*=X\theta^*$,
$\theta^*=D_x\beta^*$, $\Sz=\mathrm{supp}\,\beta^*$ and $s_0=|\Sz|$; $\Pr_0$ is the law of $\varepsilon$ and ``joint
probability'' adds the algorithm's randomness. Put $\delta_1=1-\pi(\Sz)$ and $r_n=\sqrt{\log n/n}$.
\textbf{Regime R1} is the triangular array $n\to\infty$, $s_0$ fixed, $p=p_n\le n^{A}$ for a fixed $A\ge1$.
Two design/model bundles are used.
\begin{itemize}
\item \textbf{G}: fixed design; $\varepsilon\sim N_n(0,\sigma_*^2I)$ with $\sigma_*^2>0$ fixed; intercept removed by contrasts.
\item \textbf{R}: $s_0$ fixed; $p\le n^A$ with $A\ge 1$; $\|\mu\|^2/n\le M$; $\lambda_{\min}(R_{\Sz})\ge\varphi_0^2>0$;
$s_0\le K$ eventually.
\end{itemize}

\subsection{The sampler, in the notation the results use}

% STATUS: definitions; VERIFIED-BY: n/a; SOURCE: main:Docs/Theory/theorems/00_COMMON.md:§2; main:Docs/Methodology/METHODOLOGY_v3.md (authoritative)
The method partitions $[p]$ into blocks $C_1,\dots,C_J$ (deterministic in $X$; complete linkage, top-down cut,
next-fit packing) of size at most $B=\min(\lfloor\sqrt n\rfloor,\lfloor\sqrt p\rfloor,16)$, and computes block tables
$r_b(A)=\mathrm{BF}(A)=m(A)/m(\emptyset)$ for $A\subseteq C_b$ against the full response. The \emph{composite start}
is $\pi_0=\gamma_0/\bar Z_0$ with $\gamma_0(S)=\varpi(S)\prod_br_b(S\cap C_b)$, sampled exactly; the \emph{start
mode} is $\tilde S_0=\argmax\pi_0$ (computed exactly by max-plus convolution of the block tables). A
tempering bridge $\gamma_\phi=\gamma_0e^{\phi d}$, $\phi\in[0,1]$, with $d(S)=\log\mathrm{BF}(S)-\sum_b\log r_b(S\cap C_b)$,
connects $\pi_0=q_0$ to $\pi=q_1$. One round at level $\phi$ consists of $p$ Gibbs flips, one uniform swap, $B$
neighbour swaps and one independence proposal from $\pi_0$, in a uniformly random order. The calibration
$\cC$ (grid, number of stages $T$, burn-in $L=\max(4,\lceil\log_2(n+1)\rceil)$, population size $N_e$, tier, \dots)
is a function of the data and calibration streams only. The final run uses 4 islands of $N_e$ exact i.i.d.\ draws
from $\pi_0$, the frozen grid, $L$ burn-in rounds and 4 measured rounds at $\phi=1$, and a local search of at most
$\ell=\lceil\log_2(n+p+1)\rceil$ moves whose scans score every feasible flip and swap of the (frozen) centre. Outputs
are $\Zh_r$ (island evidence), $\Zh=\tfrac14\sum_r\Zh_r$, the bound $0.05\Zh$, the Rao--Blackwell islandwise PIPs
$\PIPh_{rj}$, the primary PIP $\bar P_j=\tfrac14\sum_r\PIPh_{rj}$, $\widehat{\MPM}$ and $\Sh$. The final run scores
$Q_f:=4N_e(T+L+5)$ i.i.d.\ draws from $\pi_0$, with $4N_e(L+6)\le Q_f$.

\subsection{Runs, outputs, scored sets and LS0}

% STATUS: definition; VERIFIED-BY: W2-BOOK (single run; set inclusions only); SOURCE: W2-BOOK:STATEMENTS.tex:Potential run and output events
\begin{definition}[Potential run and output events]\label{def:run}
The \emph{potential uninterrupted run} $\bar R$ is the record the method produces under exact arithmetic with the
resource preflight and the in-run time guard removed. Let $\cG=\{\text{status}=\texttt{resource\_limit}\}$ and let
$U\supseteq\cG$ be the no-output event (with potential-run part $\bar U$). Write $\cO=U^c$ for the event that output
is available (status \texttt{exact} or \texttt{completed}), $\cO_E$ when evidence and PIPs are reported and
$\cO_H$ when an HPM is reported; $\bar{\cO}:=\bar U^c$ is the output-available event of the potential run. $\bar{\cS}$ is the scored set of $\bar R$ (calibration, final run, rejected
proposals, \emph{both} Gibbs configurations at each flip, local-search scans).
\end{definition}

% STATUS: IMPL (code reading) + method change under the investment rule (pilot data pending); VERIFIED-BY: none (HQ read of the code for centre emission); SOURCE: W4-T38FULL:T38_STATEMENT.md:Objects (LS0); main:theory_campaign/HQ/SR_AUDIT.md (decision recorded); F-REFINE:ARCHITECTURE.md:§4
\begin{definition}[LS0: the final local search from both the best-scored model and the start mode]\label{def:ls0}
The method's final local search is the production local search (at most $\ell$ moves, all feasible flips and swaps
scanned, strict-improvement moves, lexicographic tie rule) centred at the running HPM. \textbf{LS0} is the one extra
production local search centred at the start mode $\tilde S_0=\argmax\pi_0$, with its centre emitted as a score
event \emph{before} the search begins and with the same HPM tie rule. Let $\cS_{\rm LS0}$ be the models it scores and
\[
\cS^{+}:=\cS\cup\cS_{\rm LS0},\qquad \Spl:=\text{the best model of }\cS^{+}\text{ under the method's tie rule}.
\]
Without LS0, $\cS^+=\cS$ and $\Spl=\Sh$. When LS0 is part of the method (the presentation approved by the project
lead on 2 November 2026, conditional on the data check O-LS0, which was not available when this draft was written),
the reported HPM is $\Spl$, and $\bar{\cS}^+$ denotes the corresponding potential scored set.
\end{definition}

\subsection{Notation for the repair theory}

% STATUS: definitions; VERIFIED-BY: W3-VT37R, A-REV2 (verdicts on the statements that use them, below); SOURCE: W1-T37R:PROOF.md:Definitions for the refined statement; W4-T38FULL:T38_STATEMENT.md:Objects; F-REFINE:CARDS.md:Card 4
For a model $S$ put $\Phi(S)=|S\,\triangle\,\Sz|$, $m(S)=|\Sz\setminus S|$ (missing truths) and $e(S)=|S\setminus\Sz|$
(extras), so $\Phi=m+e$. A flip or swap neighbour $T$ of $S$ is a \emph{progress} neighbour if $\Phi(T)<\Phi(S)$:
adding a missing truth, deleting an extra, or swapping an extra for a truth. For $S=U\cup P$ with $U\subseteq\Sz$,
$P\cap\Sz=\emptyset$, the lattice is $\Lat(S)=\{U'\cup P':U\subseteq U'\subseteq\Sz,\ P'\subseteq P\}$, of size $2^{\Phi(S)}$.
$\cR_{a,d}=\{S=U\cup P:U\subseteq\Sz,\ m(S)\le a,\ P\cap\Sz=\emptyset,\ |P|\le d\}$ (feasible, $|S|\le K$), and
$\cR_d=\cR_{s_0,d}$. $M(S)$ is the production best-move map ($\bot$ if no strictly improving feasible flip or swap).
$D_S$ is the event that the best progress neighbour of $S$ strictly beats both $S$ and every non-progress neighbour.
$P_T$ is the orthogonal projection on $\mathrm{span}\,X_T$ and $\Delta(T)=\|(I-P_T)\mu\|^2/(n-1)$ is the
population residual (deterministic under G).

% =====================================================================================================
\section{Standing assumptions}\label{sec:assumptions}
% =====================================================================================================

The following five declared assumptions are used throughout; each theorem lists any further hypotheses itself.

% STATUS: IDEAL; VERIFIED-BY: n/a; SOURCE: main:Docs/Theory/theorems/00_COMMON.md:§3 (EX-A)
\begin{assumption}[EX-A]\label{as:exa}
Exact scores, convolutions, probabilities and comparisons (exact arithmetic).
\end{assumption}

% STATUS: IDEAL; VERIFIED-BY: n/a; SOURCE: main:Docs/Theory/theorems/00_COMMON.md:§3 (EX-S)
\begin{assumption}[EX-S]\label{as:exs}
Independent ideal uniform streams with the method's stream separation.
\end{assumption}

% STATUS: method construction (METHODOLOGY_v3 §§6.8--6.9); VERIFIED-BY: n/a; SOURCE: main:Docs/Theory/theorems/00_COMMON.md:§3 (CAL)
\begin{assumption}[CAL]\label{as:cal}
The calibration $\cC$ uses streams 0--9 only; the final streams 10--13 are fresh; escalation is decided before any
final randomness. In particular the calibration is a function of the data and of streams 0--9 only, and every
fixed-data identity below conditions on a successful frozen $\cC$, never on later diagnostics.
\end{assumption}

% STATUS: IMPL (code reading); VERIFIED-BY: HQ read of v3_engine.h step_gibbs (theory_campaign/HQ/SR_AUDIT.md; HQ's own read, not an independent audit); W3-VT38E and W2-VT24G confirmed it in code; W5-LOCKBOOK (f3c037d, unreviewed, in progress) holds a separate GUARD/OUT/SR code audit whose verdict supersedes if it differs; SOURCE: main:Docs/Theory/theorems/00_COMMON.md:§3 (SR); main:theory_campaign/HQ/SR_AUDIT.md
\begin{assumption}[SR]\label{as:sr}
Each Gibbs update evaluates both feasible configurations (the current state $S$ and the flipped state
$S\triangle\{j\}$), and both count as scored: ``scored'' means offered to the HPM store at least once in the fit.
An infeasible addition ($|S|=K$, $j\notin S$) is not evaluated.
\end{assumption}

% STATUS: IMPL (code reading of v3_runtime.h, v3_engine.h, v3_core.cpp); VERIFIED-BY: none (needs an independent code check; W5-LOCKBOOK AUDIT.md, unreviewed, covers it); SOURCE: W2-BOOK:STATEMENTS.tex:Definition (Potential run and output events); W2-BOOK:PROOF.md:§0.1
\begin{assumption}[GUARD]\label{as:guard}
The resource guard is a stop-only coupling to $\bar R$: on $\cG^c$ the status and every reported output equal those
of $\bar R$. No assumption is made on the law of $\cG$ given $\bar R$. (Operationally: no guard path touches a
random-number generator or a setting.)
\end{assumption}

% STATUS: IMPL; VERIFIED-BY: none; SOURCE: W2-T26:PROOF.md:§0 (OUT, START); W2-BOOK:PROOF.md:§1 (setting)
\begin{remark}[Two further operational assumptions]\label{rem:outstart}
\textbf{START}: the start draws are exact i.i.d.\ $\pi_0$ draws with exact $\bar Z_0$. \textbf{OUT} (used by
Theorem~\ref{thm:T26}): on $\cO\cap\{\text{SMC branch taken, calibration succeeded}\}$ the reported outputs equal the
potential run's; on $\cO\cap\{\text{exact enumeration branch}\}$ they are exact. In all statements below,
\emph{nothing is ever conditioned on $\cO$}: conclusions have the form $\Pr(\cO\cap\{\text{error}\})$.
\end{remark}

% =====================================================================================================
\section{The story}\label{sec:story}
% =====================================================================================================

% STATUS: narrative summary; every claim points to the numbered result that carries it; VERIFIED-BY: n/a (inherits the statuses of Theorems 1--6); SOURCE: F-REFINE:ARCHITECTURE.md:§1
SMC-BVS reports a consistent model because the truth gets \emph{scored} and the posterior then \emph{ranks} it first.
The reported HPM is the best model ever scored: every table and scan, both Gibbs configurations of every flip,
rejected proposals, and (with LS0) the local search from the start mode. Consequently the error decomposes, for every
$n$ and every data set, as
\begin{equation}\label{eq:story}
\Pr(\text{output and }\Spl\ne\Sz)\ \le\ 2\,\E_0[1-\pi(\Sz)]\ +\ \Pr(\text{the output-potential run never scores }\Sz)
\end{equation}
(Theorem~\ref{thm:T21}). The paper proves that the two terms vanish, in two different ways.
The \emph{posterior} term vanishes by Theorem~\ref{thm:T03} (selection consistency under anchored, size-restricted
sparse-eigenvalue conditions and beta-min). The \emph{discovery} term vanishes in one of two regimes.
\begin{itemize}
\item \textbf{Favourable regime.} The start already finds the truth: $1-\pi_0(\Sz)\to0$ (Theorem~\ref{thm:favourable}), so the final run's i.i.d.\ $\pi_0$ draws score $\Sz$ (Theorem~\ref{thm:T38}(a)).
\item \textbf{Masked regime.} The start mass at the truth vanishes, $\pi_0(\Sz)\to0$ (variables whose contribution
cancels across blocks cannot be seen by any block table), but the start concentrates \emph{near} the truth at a
reference model $S^\circ$ (Theorem~\ref{thm:masked}), and a deterministic greedy repair from the reference
scores the truth (Theorem~\ref{thm:T37}, repair), which LS0 delivers from the start mode (Theorem~\ref{thm:T38}(b1)).
\end{itemize}
Theorem~\ref{thm:T26} says that, in the favourable regime, the \emph{other} reported quantities (evidence, inclusion
probabilities) are consistent too; Proposition~\ref{prop:frozen} says that the frozen-run identities and the 95\%
evidence bound hold at every $n$ for the potential run, and at the output level under GUARD; Proposition~\ref{prop:decision}
says that the HPM is the right decision under 0--1 loss and names the one within-size error that no size-prior
retuning can fix. What the story cannot say is anything numerical at $n=200$: this is the subject of
Section~\ref{sec:disclosure}, and it must be read with every theorem.

% =====================================================================================================
\section{Results}\label{sec:results}
% =====================================================================================================

\subsection{Theorem 1: selection consistency (T03$^\sharp$) and the exact size tail}

\noindent\textbf{Notation.} Let $q_k=\binom pk w(k)$ and
\[
R_m\;=\;\frac{\varpi(\{S\in\Omega:|S|>m\})}{\varpi(\Sz)}\;=\;\sum_{k=m+1}^{K}\frac{q_k}{w(s_0)}
\]
(a data-free number; $R_m=0$ for $m\ge K$). For real $\omega$ let
$m^\sharp_n(\omega)=\min\{m\ge s_0:\ R_m\le n^{-(s_0+1)/2}e^{-\omega}\}$ and $m^\sharp_n:=m^\sharp_n(\tfrac12\log n)$
(so $s_0\le m^\sharp_n\le K$). Let $j_0=\lceil(2A+1)s_0\rceil+1$ and fix $g_0>0$.

% STATUS: definitions of the three restricted-eigenvalue/beta-min hypotheses; VERIFIED-BY: W2-VT03R (Lemma A1 for the equivalence), A-REV2 (T03♯ VERIFIED under these hypotheses); SOURCE: W1-T03R:PROOF.md:§4 (SE_λ, SE_w); §6 (BM-∞, BM(c)); W2-VT03R:REPORT.md:Lemma A1, R2
\begin{definition}[Hypotheses of Theorem~\ref{thm:T03}]\label{def:SE}
\begin{itemize}
\item $\mathrm{SE}_\lambda(j_0)$ \emph{(anchored, normalised, bounded size)}: $\lambda_{\min}(R_{\Sz\cup J})\ge\varphi_0^2$ for every
$J\subseteq[p]\setminus\Sz$ with $|J|\le j_0$ and $\Sz\cup J\ne\emptyset$.
\item $\mathrm{SE}_w(m;g_0)$ \emph{(anchored, unnormalised)}: $\lambda_{\min}(G_{\Sz\cup J})\ge g_0$ (equivalently
$\lambda_{\min}(R_{\Sz\cup J})\ge g_0/(n-1)$) for every $J\subseteq[p]\setminus\Sz$ with $|J|\le m$ and $\Sz\cup J\ne\emptyset$.
\item $\mathrm{BM}\text{-}\infty$: $\min_{j\in\Sz}\theta_j^{*2}\ge M_n\sigma_*^2\log n/(n\varphi_0^2)$ with $M_n\to\infty$.
$\mathrm{BM}(c)$: the same with a fixed constant $c>c_0(s_0,A)$, where $c_0(s_0,A)<\infty$ depends only on $s_0$ and
$A$ (the proof gives $c_0=L_*=\max\{4C_*,32[s_0(A+1)+2]\}$; this is a proof artefact, not sharp).
\end{itemize}
By eigenvalue interlacing, ``anchored'' is equivalent to the unanchored condition on every column set $T$ with at
most $j_0$ (resp.\ $m$) columns outside $\Sz$; in particular $\mathrm{SE}_w(m;g_0)$ requires every $m$ columns of $X$
to be linearly independent, with $\lambda_{\min}(X^\top X)\ge g_0$ on them.
\end{definition}

% STATUS: PROVED (candidate); VERIFIED-BY: W2-VT03R (no GAP/ERROR in the mathematics), A-REV2 (T02(i♯) VERIFIED: "explicit constant and pointwise-in-m uniform tail-probability bound are correct"; E1, E2, E3 VERIFIED); SOURCE: W1-T03R:PROOF.md:Theorem T02(i♯), Lemma E1, E2
\begin{lemma}[Exact size tail, T02$(\mathrm{i}^\sharp)$]\label{thm:T02}
Assume G and R, and let $n$ be such that $s_0\le K$. For every integer $m$ define $\Theta^\sharp_{n,m}\ge0$ by
$\Pi(|S|>m\mid y)=R_m\,(n+1)^{s_0/2}\sqrt n\;\Theta^\sharp_{n,m}$ if $R_m>0$ and $\Theta^\sharp_{n,m}:=0$ if $R_m=0$.
Then for every $L>c^{\sharp2}$ (with $c^\sharp$ and $C_\Lambda$ explicit constants depending on $\sigma_*^2,M,\varphi_0$ only)
\[
\sup_m\Pr_0(\Theta^\sharp_{n,m}>L)\ \le\ \frac{8+24M/\sigma_*^2}{n-1}+\frac1{\sqrt L}+\frac{C_\Lambda}{\log(\sqrt L/c^\sharp)} .
\]
Consequently, for every deterministic real sequence $\omega_n$, $\Pi(|S|>m^\sharp_n(\omega_n)\mid y)=O_{\Pr_0}(e^{-\omega_n})$; in
particular $\Pi(|S|>m^\sharp_n\mid y)=O_{\Pr_0}(n^{-1/2})$, with $m^\sharp_n\le\min\{K,m_n^{\rm old}\}$. The inequality is
non-asymptotic. No sparse-eigenvalue condition is used.
\end{lemma}
\noindent\emph{Proof:} see Supplement \S S.1 (source branch \texttt{theory/W1-T03R}, \texttt{PROOF.md} \S\S1--2).

% STATUS: PROVED (candidate); VERIFIED-BY: W2-VT03R, A-REV2 ("T02(ii♯)... Fixed unnormalized g0>0 suffices on the anchored comparison unions": VERIFIED); SOURCE: W1-T03R:PROOF.md:§5 Theorem T02(ii♯)
\begin{lemma}[Superset mass, T02$(\mathrm{ii}^\sharp)$]\label{thm:T02ii}
In regime R1 under G, R and $\mathrm{SE}_w(m^\sharp_n;g_0)$: $\Pi(\{S\in\Omega:S\supsetneq\Sz\}\mid y)=O_{\Pr_0}(r_n)$.
\end{lemma}
\noindent\emph{Proof:} see Supplement \S S.1 (source branch \texttt{theory/W1-T03R}, \texttt{PROOF.md} \S5).

% STATUS: PROVED (candidate) as T03♯; the lines "no beta-min", "scale invariance" are conclusions (iv),(v) of the same statement; VERIFIED-BY: W2-VT03R (with threshold correction R1 and wording R2), A-REV2 ("T03♯(i)-(v), under its SE and BM hypotheses": VERIFIED; "no hidden normalized SE use on large unions found"); SOURCE: W1-T03R:PROOF.md:§6 Theorem T03♯ (hypotheses and conclusions copied); F-REFINE:CARDS.md:Card 1
\begin{theorem}[Selection consistency, T03$^\sharp$]\label{thm:T03}
Assume regime R1, G, R, $\mathrm{SE}_\lambda(j_0)$, $\mathrm{SE}_w(m^\sharp_n;g_0)$ for a fixed $g_0>0$, and a beta-min condition
in one of the forms $\mathrm{BM}\text{-}\infty$ or $\mathrm{BM}(c)$ of Definition~\ref{def:SE} (for (i)--(iii)). No hypothesis H-T02 is used
(Theorem~\ref{thm:T02} is used directly). Then:
\begin{enumerate}
\item[(i)] $1-\pi(\Sz)=O_{\Pr_0}(r_n)$;
\item[(ii)] $\max_j|\PIP_j-\Ind\{j\in\Sz\}|\le1-\pi(\Sz)$ for \emph{every} data set, and $\Pr_0(\MPM=\Sz)\to1$;
\item[(iii)] $\Pr_0(S^\dagger\ne\Sz)\le2\E_0[1-\pi(\Sz)]\to0$, where $S^\dagger$ is the global HPM; on $\{\pi(\Sz)>\tfrac12\}$,
$\Sz$ is the strict global maximiser of $F$ (so $\HPM=\MPM=\Sz$ strictly);
\item[(iv)] (no beta-min) for fixed $c>c_0$ and $J_n(c)=\{j\in\Sz:\theta_j^{*2}\ge c\,\sigma_*^2\log n/(n\varphi_0^2)\}$:
$\Pi(S\not\supseteq J_n(c)\mid y)=O_{\Pr_0}(r_n)$ and $1-\min_{j\in J_n(c)}\PIP_j=O_{\Pr_0}(r_n)$; the same holds for
$J_n^\infty=\{j:\theta^{*2}_j\ge M_n\sigma_*^2\log n/(n\varphi_0^2)\}$ with any $M_n\to\infty$;
\item[(v)] scale invariance of the conclusions.
\end{enumerate}
\end{theorem}
\noindent\emph{Proof:} see Supplement \S S.1 (source branch \texttt{theory/W1-T03R}, \texttt{PROOF.md} \S6; locked T01--T03 as cited there).

% STATUS: PROVED (W2-VT03R Lemma A1); PROOF SKETCH (Prop. W); VERIFIED-BY: W2-VT03R (A1: proved there); A-REV2 did not review Prop. W; SOURCE: W2-VT03R:REPORT.md:Lemma A1, Proposition W
\begin{remark}[Scope of Theorem~\ref{thm:T03}]\label{rem:T03}
\begin{enumerate}
\item The ``anchoring'' of the eigenvalue condition is cosmetic (interlacing; Definition~\ref{def:SE}). The real
relaxations relative to the locked SE are the size range ($m^\sharp_n$ instead of $m_n$) and the normalisation ($g_0$
instead of $(n-1)\varphi_0^2$) on large unions.
\item The statement ``the exact tail cuts $m_n$ to $s_0+O(s_0\log n/\log p)$'' is \emph{false} under G and R alone
(duplicated columns, W1-T03R Prop.~D, COUNTEREXAMPLE, reproduced by W2-VT03R and A-REV2); the exact tail improves the
locked size cutoff by a constant ($4s_0!\,e^{\gamma_m}$) and an $m^2/(4p)$ term, not by a power of $p$. [Status: PROVED/VERIFIED, W1-T03R E1(d), Prop.~D.]
\item (PROOF SKETCH, W2-VT03R Prop.~W.) $\mathrm{SE}_\lambda(j_0)+\mathrm{BM}$ may be replaced by $\mathrm{SE}_w$ plus a direct
signal-separation condition SIG$(c)$, $c\ge L_*$, for conclusions (i)--(iii).
\end{enumerate}
\end{remark}

\subsection{Proposition 1: Bayes decisions and the size-prior ceiling (T07, with T06(i))}

% STATUS: PROVED-LOCKED (T07(i)--(iii), T06(i)(a)) with W2-BOOK restatement (i′),(iii′) PROVED (single run); VERIFIED-BY: none beyond the locked review chain; W5-LOCKBOOK (f3c037d) has this on its task list but pushed no report; SOURCE: W2-BOOK:STATEMENTS.tex:Proposition (Decision rules and the size-prior ceiling) [label prop:decision]; W2-BOOK:PROOF.md:§3; F-REFINE:CARDS.md:Card 2
\begin{proposition}[Decision rules and the size-prior ceiling]\label{prop:decision}
Let $\pi$ be the target posterior and $S^*_k$ the tie-order-least maximiser of the marginal likelihood among
models of size $k$.
\begin{enumerate}
\item[(i)] The HPM is a Bayes rule for $0$--$1$ loss on supports; the median probability model is a Bayes rule for Hamming loss.
\item[(ii)] For any PIP estimate $\hat p$ the plug-in Hamming rule has excess posterior risk at most $2\|\hat p-\PIP\|_1$.
\item[(iii)] If $s_0\le K$, some positive size-only prior has global HPM $\Sz$ if and only if $S^*_{s_0}=\Sz$.
In particular, if a scored model $S'\ne\Sz$ with $|S'|=s_0$ has a higher score than $\Sz$, no positive size-only prior
recovers $\Sz$ on those data.
\item[(iv)] Two supports have equal score for every response if and only if $X_SX_S^\top=X_{S'}X_{S'}^\top$ (outer
products, not projections); they then have equal size.
\end{enumerate}
\end{proposition}
\noindent\emph{Proof:} see Supplement \S S.2 (source branch \texttt{theory/W2-BOOK}, \texttt{PROOF.md} \S3; locked T07, T06).

% STATUS: wording constraint (not a theorem); VERIFIED-BY: n/a; SOURCE: F-REFINE:CARDS.md:Card 2 (Risk); F-REFINE:ARCHITECTURE.md:§5
\begin{remark}\label{rem:ceiling}
Part (iii) concerns \emph{within-size} errors ($|S|=s_0$, $m(S)\ge m(\Sz)$) only; size-prior retuning \emph{can} move
different-size (superset or subset) errors. It covers every truncated BB$(a,b)$ prior. For (iv), with the cap $K$ the
MPM keeps the $\min(K,N_+)$ largest PIPs above $1/2$ and is unique iff no $\PIP_j=\tfrac12$.
\end{remark}

\subsection{Theorem 2: the start --- favourable and masked regimes (T22, T27)}
\stepcounter{theorem}

\noindent\textbf{Setting.} G, R-lite$(A,M)$ (meaning: $p\le n^A$ and $\|\mu\|^2\le Mn$), the v3 partition, blocks of size
at most $16$, $\nu=0.01+n/2$. \textbf{SE-blk}$(\varphi_0)$: $\min_b\lambda_{\min}(R_{C_b})\ge\varphi_0^2$ (checkable on the realised
blocks; by interlacing equivalent to $\lambda_{\min}(X_U^\top X_U/(n-1))\ge\varphi_0^2$ for every nonempty $U$ inside a block).
\textbf{BO}: $X_{C_b}^\top X_{C_{b'}}=0$ for $b\ne b'$. For a reference $S^\circ$ and block $C_b$, $U_b^\circ=S^\circ\cap C_b$,
$\lambda_b(S^\circ)=\|(P_{C_b}-P_{U_b^\circ})\mu\|^2$ and $g_b(S^\circ)=\min_{U\subseteq C_b,U_b^\circ\not\subseteq U}\|(P_{U\cup U_b^\circ}-P_U)\mu\|^2$.
The conditions VIS, $\mathrm{NL}_\eta$, $\mathrm{NL}_{\rm avg}(\kappa)$, EXP$'$ are those of Section~3 of the common-definitions file.
Specific conditions used below: \textbf{ZL}$(S^\circ)$ (exact zero leakage): $\lambda_b(S^\circ)=0$ for all blocks;
\textbf{SPREAD}$(\Delta_0)$: $\min_b\|(I-P_{C_b})\mu\|^2/(n-1)\ge\Delta_0>0$; \textbf{VIS-lin}$(c_V)$: $\min_bg_b(S^\circ)\ge c_Vn$ with fixed
$c_V>0$; \textbf{BM-blk}: $\min_{j\in\Sz}\theta_j^{*2}\ge M_n(\sigma^2+2M)\log n/(\varphi_0^2(n-1))$ with $M_n\to\infty$;
\textbf{CL}$(L)$ (chain leakage): along a deterministic addition chain $S^{(0)}\subset\cdots\subset S^{(c)}=\Sz$ with
$S^{(i+1)}=S^{(i)}+l_i$, $\max_i|\delta_{T_i,l_i}|\le L\sigma$, where $T_i=S^{(i)}\cap C_{b(l_i)}$ and
$\delta_{T,l}=u_{T,l}^\top\mu$ ($u_{T,l}$ the unit residual of $x_l$ on $X_T$); \textbf{H3}: $\Pr_0(\pi_0(S^{(0)})\ge c_0)\to1$.

% STATUS: PROVED-LOCKED (T22(i),(ii),(iv): wave 1, full chain); VERIFIED-BY: locked review chain; W2-T26 Lemma 1 (the Δ_n identity); SOURCE: F-REFINE:CARDS.md:Card 3 (i); W2-T26:PROOF.md:Lemma 1
\begin{thmpart}[deterministic start bounds]\label{thm:T22}
\emph{(i)} The composite start puts mass at least $(1-q)^{s+1}/(1+e)^s$ on any feasible reference model $S^\circ$
($S^\circ\ne\Sz$ allowed), with $q,e$ the table constants of the locked T22(i); the odds from a reference to $\Sz$ are exact;
and the tempering bridge keeps $\sup_{\phi\in[0,1]}[1-q_\phi(\Sz)]\le\max\{1-\pi_0(\Sz),1-\pi(\Sz)\}$. Hence, if $\Sz\in\Omega$,
$\Delta_n:=1-\inf_{\phi}q_\phi(\Sz)=\max\{1-\pi_0(\Sz),\,1-\pi(\Sz)\}<1$ exactly.
\end{thmpart}
\noindent\emph{Proof:} see Supplement \S S.3 (locked T22; source branch \texttt{theory/W2-T26}, \texttt{PROOF.md} Lemma 1).

% STATUS: PROVED; VERIFIED-BY: W3-VEXP (EXP′ theorem and finite-n certificate R1), A-REV2 ("SEX Theorem 1 EXP′, exact remainder and stated bound: VERIFIED" under fixed-X actual blocks, uniform SE-blk, bounded B, polynomial family; the displayed limiting constant 0.52/c_* is an ERROR, the actual limit is 0.51/c_*, so it is not quoted here); SOURCE: W1-EXP:PROOF.md:Theorem 1 (EXP′); W3-VEXP:REPORT.md:R1
\begin{thmpart}[EXP$'$ is a theorem]\label{thm:EXP}
Assume G, R-lite$(A,M)$ and SE-blk$(\varphi_0)$. Put $c_*=\sigma^2/(8(M+2\sigma^2))$. There is an explicit $\kappa_{\rm EXP}$
(depending on $\varphi_0,\sigma,M,n$ only) and an event $E_n$ with $\Pr_0(E_n)\to1$ such that on $E_n$, \emph{every}
within-block addition $(b,T,l)$ ($T\subseteq C_b\setminus\{l\}$) satisfies
\[
\log\frac{r_b(T+l)}{r_b(T)}=-\tfrac12\log(n+1)-\nu\log(1-x_{T,l})+R_{n,b,T,l},\qquad |R_{n,b,T,l}|\le\kappa_{\rm EXP},
\]
where $x_{T,l}=(\mathrm{RSS}_T-\mathrm{RSS}_{T+l})/\mathrm{RSS}_T$. In particular EXP$'$ holds in its strict form ($\rho_n=\max|R|=O_{\Pr_0}(1)$);
at finite $n$ a data-computable certificate replaces $\kappa_{\rm EXP}$. Consequently EXP$'$ is removed from the hypothesis list of the
favourable-start statement, and $\mathrm{NL}_\eta$ may be weakened to $\mathrm{NL}_C$ (W1-EXP Prop.~2.1, PROOF SKETCH).
\end{thmpart}
\noindent\emph{Proof:} see Supplement \S S.3 (source branch \texttt{theory/W1-EXP}, \texttt{PROOF.md} \S1; \texttt{theory/W3-VEXP}).

% STATUS: (F1) PROVED-LOCKED CONDITIONAL on EXP′ (now discharged) with W1-EXP Prop 2.1 PROOF SKETCH; (F2) PROOF SKETCH in W1-EXP (verbatim corollary of (F1)+Thm 1); (F3) PROVED; VERIFIED-BY: W3-VEXP ((F2) VERIFIED as a verbatim corollary, R2; (F3) VERIFIED); SOURCE: F-REFINE:CARDS.md:Card 3 (iii); W1-EXP:PROOF.md:Prop 2.1, Theorem 2, Theorem 6
\begin{thmpart}[favourable start]\label{thm:favourable}
For a feasible reference $S^\circ$ (which may be $\Sz$, or $S^\circ\ne\Sz$), $1-\pi_0(S^\circ)\to0$ in probability at the rate
stated under any one of:
\begin{itemize}
\item[(F1)] the locked T22(iii) hypotheses minus EXP$'$: G, R, SE-blk, VIS, $\mathrm{NL}_\eta$, $\mathrm{NL}_{\rm avg}(\kappa)$ (all at reference $S^\circ$): rate $O_{\Pr_0}(2^B\sqrt{\log n}\,n^{-\kappa})$;
\item[(F2)] BO and BM-blk (with G, R, SE-blk, reference $\Sz$): rate $O_{\Pr_0}(2^B\sqrt{\log n}\,n^{-1/2})$;
\item[(F3)] G, R-lite, SE-blk, ZL$(S^\circ)$, SPREAD, VIS-lin$(S^\circ)$: rate $O_{\Pr_0}(2^B n^{-1/2})$, self-contained (no locked-proof transfer).
\end{itemize}
\end{thmpart}
\noindent\emph{Proof:} see Supplement \S S.3 (source branch \texttt{theory/W1-EXP}, \texttt{PROOF.md} \S2, Theorem 6; verification \texttt{theory/W3-VEXP}).

% STATUS: PROVED (Theorems 3 and 4 of W1-EXP given H3); the limit-law refinement is a PROOF SKETCH; VERIFIED-BY: W3-VEXP; A-REV2 ("SEX Theorem 3 deterministic fixed-length chain: VERIFIED"; "SEX Theorem 4, fixed CL(L) and normalized H3: VERIFIED"; "growing-leakage bound: VERIFIED as a bound, the 'iff' gloss is an ERROR, sufficiency only"); SOURCE: W1-EXP:PROOF.md:Theorems 3, 4 and Limit law; F-REFINE:CARDS.md:Card 3 (iv)
\begin{thmpart}[masked start]\label{thm:masked}
Assume G, R-lite, SE-chain (implied by SE-blk), a deterministic fixed-length addition chain, CL$(L)$ with $L$ fixed, H3 and
$p\ge2(s+c)$ ($s=|S^{(0)}|$, $c\ge1$ the chain length). Then with probability $\to1$, for all $i<c$,
\[
-1-\kappa_{\rm EXP}\ \le\ \log\frac{r_{b_i}(T_i+l_i)}{r_{b_i}(T_i)}+\tfrac12\log n\ \le\ \kappa_{\rm EXP}+1.01(L+|Z_i|)^2\quad(n\ge400),
\]
with $Z_i\sim N(0,1)$, and
\[
\log\pi_0(\Sz)+c\log p+\tfrac c2\log n=O_{\Pr_0}(1),\qquad\text{i.e. }\ \pi_0(\Sz)=p^{-c}n^{-c/2}e^{O_P(1)} .
\]
With $L=L_n\to\infty$ the upper side gains the term $1.01c(L_n+O_P(1))^2$ (a sufficient bound only). This is a statement about
i.i.d.\ draws from $\pi_0$, \emph{not} about the hybrid sampler.
\end{thmpart}
\noindent\emph{Proof:} see Supplement \S S.3 (source branch \texttt{theory/W1-EXP}, \texttt{PROOF.md} \S3).

% STATUS: exact-Gram part PROVED (W1-EXP Theorem 5 (a),(b),(c),(e); (d) by locked T03; the limit law inside (b) is a PROOF SKETCH); iid-Gaussian p≫n part PROVED by one run, unverified; VERIFIED-BY: W3-VEXP (exact Gram; |ψ|<ρ boundary sharp in effect, row 17); iid part: none (W3-PXN single run; W5-VMASK, branch c3e1b73, has pushed an independent core and sanity checks only, no verdict); SOURCE: W1-EXP:PROOF.md:Theorem 5, Lemmas 4.1, 4.2; W3-PXN:PROOF.md:P1–P8; F-REFINE:CARDS.md:Card 3 (v)
\begin{thmpart}[one family carries both]\label{thm:LH}
\textbf{The family LH} (linked hubs): $G$ groups of size $g$ with equicorrelation $\rho_h\in[\rho_-,\rho_+]\subset(0,1)$, a hub in each group;
groups independent except $c$ disjoint \emph{linked pairs} $(a_i,b_i,\psi_i)$ with cross-covariance through the hubs; truth
$\theta_{j_{a_i}}=\beta_i\ne0$, $\theta_{j_{b_i}}=-\psi_i\beta_i$ (the \emph{suppressor}), further unlinked visible groups with fixed
nonzero coefficients; $s_0=2c+\#\text{visible}$ fixed. The masked set is $M=\{j_{b_1},\dots,j_{b_c}\}$ and $S^\circ=\Sz\setminus M$.
\emph{Exact-Gram realisation:} $X^\top X=(n-1)\Sigma_{\rm LH}$, $p\le n-1$, $|\psi_i|<\min(\rho_{a_i},\rho_{b_i})$, $\psi_i\ne0$,
$g\le B_n<2g$ for all large $n$, fixed coefficients. Then, under G with $\sigma^2>0$ fixed:
\begin{enumerate}
\item[(a)] $1-\pi_0(S^\circ)=O_{\Pr_0}(2^Bn^{-1/2})$, hence $\pi_0(\cR_d)\to1$ for every $d\ge0$;
\item[(b)] $\pi_0(\Sz)=p^{-c}n^{-c/2}e^{O_P(1)}$ (masking), with an exact limit law for the masking statistic (PROOF SKETCH);
\item[(c)] VIS fails at reference $\Sz$, so the favourable-start statement does not apply;
\item[(d)] $1-\pi(\Sz)=O_{\Pr_0}(\sqrt{\log n/n})$ (locked T03);
\item[(e)] $\pi_0(\Sz)\to0$, $\pi(\Sz)\to1$, $\pi_0(\cR_0)\to1$: the start is masked, the target is not.
\end{enumerate}
\emph{iid-Gaussian realisation (single unverified run).} With i.i.d.\ Gaussian rows and $n^a\le p\le n^A$ ($p\gg n$ allowed),
the partition equals the groups with probability $\to1$, $1-\pi_0(S^\circ)=O_{\Pr_0}(\sqrt{\log n/n})$ \emph{without} the $2^B$ factor,
and the exact limit law of the masking statistic holds (W3-PXN P1--P4, with P5--P8).
\end{thmpart}
\noindent\emph{Proof:} see Supplement \S S.3 (source branch \texttt{theory/W1-EXP}, \texttt{PROOF.md} \S4; \texttt{theory/W3-PXN}, \texttt{PROOF.md}).

% STATUS: wording constraint; VERIFIED-BY: n/a; SOURCE: F-REFINE:CARDS.md:Card 3 (Regime discipline, Risk); W3-PXN:PROOF.md:Prop. Q
\begin{remark}[Discipline]\label{rem:regime}
Parts (iii) and (iv)--(v) of Theorem~2 are never combined on one sequence; on LH the reference in (iii) is $S^\circ$. The cliff
$|\psi|<\rho$ (beyond it mass moves to a different regime) and the packing condition $g\le B<2g$ are part of the theorem and must be
stated. The rates in (iii) carry a factor $2^B$ and are numerically vacuous below $n\approx10^9$ (Section~\ref{sec:disclosure}).
\end{remark}

\subsection{Theorem 3: the repair (refined T37, T24)}
\stepcounter{theorem}

\noindent\textbf{Conditions.} For a finite family $\cR\ni\Sz$ of feasible models: \textbf{(L)}: $M(\Sz)=\bot$;
$\mathrm{TD}'_\infty(\cR)$: (L) and for every $S\in\cR\setminus\{\Sz\}$, $M(S)\ne\bot$ and $M(S)\in\cR$;
$\mathrm{TD}'_P(\cR)$: (L) and $M(S)$ is a progress move for every $S\in\cR\setminus\{\Sz\}$.
(Strict separation $\bigcap_SD_S$ together with (L) is \emph{sufficient} for $\mathrm{TD}'_P$; with exact ties it is not necessary.)
$\mathrm{POS}(\cR_d)$: at every $S\in\cR_d\setminus\{\Sz\}$ the move designated by literal TD(1) or TD(2) has gain $>0$.

% STATUS: PROVED (negative result); VERIFIED-BY: W3-VT37R (exhaustive), A-REV2 ("S37 Proposition 1 and swap identity: VERIFIED"); SOURCE: W1-T37R:PROOF.md:Proposition 1; W3-VT38E:V1; F-REFINE:CARDS.md:Card 4 ("What is false and must go")
\begin{lemma}[The literal greedy-descent condition is unsatisfiable]\label{prop:literalTD}
For $d\ge1$ the literal condition $\mathrm{TD}(d,m)$ of the original T37 is unsatisfiable: at $S=\Sz-j+l$ the swap $l\to j$ beats
the addition of $j$ whenever $F(\Sz)>F(\Sz+l)$, so $\mathrm{TD}(d,0)$ fails for every $d\ge1$ on $\{\pi(\Sz)>\tfrac12\}$, i.e.\ under the conclusion of
Theorem~\ref{thm:T03}. Therefore the literal T37, T37(ii)'s union bound (it includes a comparison with the wrong sign) and
``exactly $|S\triangle\Sz|$ moves'' (swaps cut $\Phi$ by $2$) are withdrawn.
\end{lemma}
\noindent\emph{Proof:} see Supplement \S S.4 (source branches \texttt{theory/W1-T37R} Prop.~1; \texttt{theory/W3-VT38E} V1).

% STATUS: PROVED; VERIFIED-BY: W3-VT37R (exhaustive check, 720 designs × 632 states at p=15, 0 violations); A-REV2 ((A1)–(A4): VERIFIED; (A5) literal TD ⇒ TD′_P without POS: ERROR, the POS repair R1 VERIFIED; "TD′_P iff (L) and all D_S": ERROR (sufficient only); "weakest" prose: ERROR (use W3-VT37R R3's correction)); SOURCE: W1-T37R:PROOF.md:Theorem A (A1)–(A5); W3-VT37R:REPAIRS.md:R1, R3; F-REFINE:CARDS.md:Card 4 part A
\begin{thmpart}[deterministic greedy repair --- Theorem A]\label{thm:T37}
Assume EX-A; $X,y$ fixed; $\ell$ the production move cap; $\cR$ a finite family of feasible models containing $\Sz$.
\begin{enumerate}
\item[(A1)] $\mathrm{TD}'_\infty(\cR)$ holds if and only if, for every $S\in\cR$, the uncapped greedy path from $S$ stays in $\cR$ and terminates at $\Sz$.
\item[(A2)] Under $\mathrm{TD}'_P(\cR_{a,d})$, from every $S\in\cR_{a,d}$ the greedy path stays in $\Lat(S)\subseteq\cR_{a,d}$, $\Phi$ strictly
decreases at each move and $\Sz$ is reached after $\tau(S)=m(S)+e(S)-\sigma(S)\le\Phi(S)$ moves, where $\sigma(S)$ is the number of
extra$\to$truth swaps used; $\max\{m(S),e(S)\}\le\tau(S)$.
\item[(A3)] If $\mathrm{TD}'_P(\cR_{a,d})$ holds and $\ell\ge a+d$ (or merely $\ell\ge\tau(S_c)$), then from any centre $S_c\in\cR_{a,d}$ the production
local search returns $\Sz$ with \texttt{local\_certified}, without \texttt{local\_limit}, and afterwards $\Sh=\Sz$ up to exact ties
(with \emph{no} $\pi(\Sz)>\tfrac12$ premise). The same holds for one centre $S_c$ under the lattice condition
(L)\,$\wedge\bigcap_{S'\in\Lat(S_c)\setminus\{\Sz\}}D_{S'}$, which involves $2^{\Phi(S_c)}$ states only.
\item[(A5, repaired)] Literal strict $\mathrm{TD}(d,0)\wedge\mathrm{POS}(\cR_d)\Rightarrow\mathrm{TD}'_P(\cR_d)\Rightarrow\mathrm{TD}'_\infty(\cR_d)$; the converse of the
second implication fails.
\end{enumerate}
\end{thmpart}
\noindent\emph{Proof:} see Supplement \S S.4 (source branch \texttt{theory/W1-T37R}, \texttt{PROOF.md} Theorem A; repairs \texttt{theory/W3-VT37R}, \texttt{REPAIRS.md} R1, R3).

% STATUS: PROVED; VERIFIED-BY: same-family referee inside W4-T38FULL (not cross-family; A-REV1/A-REV2 did not review W4-T38FULL); SOURCE: W4-T38FULL:T38_PROOF.md:§3.1 Lemma LS-score; F-REFINE:CARDS.md:Card 4 part A (LS-score)
\begin{lemma}[Lemma LS-score]\label{lem:LSscore}
If $D_S$ holds on $\Lat(C_0)\cap\{\Phi\ge2\}$ and $\ell\ge\Phi(C_0)-1$, the production search from $C_0$ \emph{scores} $\Sz$ (the next scan at
a flip neighbour scores it); certification is not required. This is the form used by Theorem~\ref{thm:T38}.
\end{lemma}
\noindent\emph{Proof:} see Supplement \S S.4 (source branch \texttt{theory/W4-T38FULL}, \texttt{T38\_PROOF.md} \S3.1).

% STATUS: PROVED for the inclusion (new-delivery monotonicity); the posterior-miss inference of the source's Lemma 0(b) is an ERROR (A-REV2) and is NOT stated; VERIFIED-BY: A-REV2 ("S37 Lemma 0(a),(c): VERIFIED"; "Lemma 0(b) as written: ERROR — new-delivery monotonicity inclusion survives; its post-search 'posterior miss' inference does not"); W3-VT37R R2 (the proposed exhaustive failure dichotomy is itself an ERROR per A-REV2); SOURCE: W1-T37R:PROOF.md:Lemma 0; W3-VT37R:REPAIRS.md:R2; A-REV2:REVIEW.md:§4.2
\begin{lemma}[Lemma 0, restricted to what survives review]\label{lem:lemma0}
The final local search, centred at the running HPM, can deliver $\Sz$ (i.e.\ $\Sz\notin\cS_{\rm pre}$ and $\Sz\in\cS$ after the search)
only on $E_{\rm LS}=\{\Sz\notin\cS_{\rm pre},\ F(\Sz)>F(\text{centre})\}$, by monotonicity of $F$ along the greedy path.
After a certified repair with $\Sz$ the centre, $\Sh=\Sz$ up to exact ties.
\end{lemma}
\noindent\emph{Proof:} see Supplement \S S.4 (source branch \texttt{theory/W1-T37R}, \texttt{PROOF.md} Lemma 0 (a),(c) only).

% STATUS: PROVED (candidate); VERIFIED-BY: none final (REPORT of W3-PGD is a stub; W5-VPGD, branch 194e12c, has pushed a same-family referee file REFEREE.md labelled "draft verdicts": no ERROR in any proof step, four minor GAPs G1-G4 and one misleading interpretive remark; not a final report, status not upgraded); A-REV2 flags the neighbouring W1-T37R Theorem C extensions as GAP and promotes no asymptotic TD′ theorem; SOURCE: W3-PGD:PROOF.md:Theorem W, Definition LD♭, Lemma S; F-REFINE:CARDS.md:Card 4 part B
\begin{thmpart}[asymptotic repair --- Theorem W]\label{thm:W}
Assume regime R1; $s_0$, $a\le s_0$, $d$ fixed; $\mathcal D=\mathcal D_n\subseteq\Sz^c$ deterministic ($X$-measurable); EX-A, G, R,
$\mathrm{SE}_\lambda(d+1)$ ($\lambda_{\min}(R_{\Sz\cup J})\ge\varphi_0^2$ for $|J|\le d+1$), $\mathrm{BM}\text{-}\infty$, \textbf{PGD}$(\eta)$ on
$\cR_{a,d}(\mathcal D)=\{U\cup P:U\subseteq\Sz,\ |\Sz\setminus U|\le a,\ P\subseteq\mathcal D,\ |P|\le d\}$ for a fixed $\eta\in(0,1]$ (at every state missing a truth,
the best non-progress reduction of the population residual is at most $(1-\eta)$ times the best progress reduction), and one of
(H1) $d\le1$; (H2) $(|\mathcal D|+1)^{d-1}\le n^{1/2-\epsilon}$; (H3) $\mathrm{LD}^\flat(d;B',\rho^*)$ with $\tau^2<1+1/(2A)$,
$\tau=(1+\rho^*\sqrt d)/\sqrt{1-\rho^{*2}}$. Then $\Pr_0\big(\mathrm{TD}'_P(\cR_{a,d}(\mathcal D))\big)\to1$; consequently for every
(data-dependent) centre $S_c$ and all $n$ with $\ell_n\ge a+d$,
\[
\Pr\big(\text{the search from }S_c\text{ does not return }\Sz\text{ certified}\big)\ \le\ \Pr\big(S_c\notin\cR_{a,d}(\mathcal D)\big)+o(1).
\]
\end{thmpart}
\noindent\emph{Proof:} see Supplement \S S.4 (source branch \texttt{theory/W3-PGD}, \texttt{PROOF.md} \S4). The condition $\mathrm{LD}^\flat$ is that
of \texttt{PROOF.md} \S3 (bounded neighbour sets $\mathfrak N(l)$ of size $\le B'$ and a squared partial multiple correlation at most $\rho^*$).

% STATUS: PROVED (W3-PGD P4: sharpness; C-LD: counterexample); VERIFIED-BY: none; W1-T37R C3: COUNTEREXAMPLE VERIFIED by W3-VT37R (PGD is not implied by SE + BM-∞); SOURCE: W3-PGD:PROOF.md:Prop P4, Prop C-LD; W1-T37R:PROOF.md:Theorem C (C3)
\begin{remark}[Limits of Theorem~\ref{thm:W}]\label{rem:Wlimits}
PGD is not implied by $\mathrm{SE}+\mathrm{BM}\text{-}\infty$ (combination look-alike counterexample; W1-T37R, verified by W3-VT37R), and is sharp: a
reverse margin $\rho^0\ge(1+\epsilon)\rho^+$ on a subsequence forces failure (W3-PGD Prop.~P4). The local-dependence condition cannot be dropped
for $d\ge2$ with unrestricted extras (equicorrelated nulls with $r>\tfrac12$: $\mathrm{TD}'_P(\cR_{0,2})$ fails with probability $\to1$; W3-PGD Prop.~C-LD,
COUNTEREXAMPLE); it is unnecessary for $d=1$ or for extras from a class of size $n^{o(1)}$.
\end{remark}

% STATUS: PROVED; VERIFIED-BY: same-family referee inside W4-T38FULL, with a brute-force population check of the margins (scripts/ref_lhlat.txt, ref_lhformulas.txt); not cross-family; SOURCE: W4-T38FULL:T38_PROOF.md:§3.2 Lemma M, §3.4 Prop. LH-LAT; F-REFINE:CARDS.md:Card 4 part C
\begin{thmpart}[explicit margin on the masked family LH]\label{prop:LHLAT}
On LH (exact Gram, as in Theorem~\ref{thm:LH}) with $S^\circ=\Sz\setminus M$, the population condition $\mathrm{POP\text{-}LAT}(S^\circ,m_0)$ holds with the explicit
sharp margin $m_0=\min\{(1-\kappa_{\max})R_{\min},\ \min_i\beta_i^2(1-\psi_i^2)^2,\ \min_h\theta_h^2\}$, where
$R_i=\beta_i^2\psi_i^2(1-\psi_i^2)$ and $\kappa_i=\rho^2(1-\psi_i^2)/(1-\psi_i^2\rho^2)$. Moreover (Lemma M) any population margin $m_n\gg\sqrt{\log n/n}$
turns $\mathrm{POP\text{-}LAT}$ into the event $D_S$ on $\Lat(S^\circ)$ with probability $\to1$; hence the search from $S^\circ$ scores $\Sz$ in $c$ moves.
\end{thmpart}
\noindent\emph{Proof:} see Supplement \S S.4 (source branch \texttt{theory/W4-T38FULL}, \texttt{T38\_PROOF.md} \S\S3.2, 3.4).

% STATUS: PROVED (finite inequality, after the Lemma F′ repair); Proposition U (the potential-function bookkeeping) is an ERROR as written and is NOT cited; the asymptotic premises are PROVED-UNDER / PROOF SKETCH / FALSE-without-conditions as listed; VERIFIED-BY: W2-VT24G (after Lemma F′); A-REV1 ("Theorem B, restated T24, finite inequality: VERIFIED after F′ repair"; Corollary P1: VERIFIED; Cor B-M: VERIFIED under IND-T; IND-T: VERIFIED for the uninterrupted process; Prop. U: ERROR; "W1 §6 premises imply Hη→0": GAP); SOURCE: W1-T24G:PROOF.md:§4 Theorem B; W2-VT24G:PROOF.md:Lemma F′, B1′/B2′; W2-PREM:PROOF.md:IND-T; F-REFINE:CARDS.md:Card 4 part D
\begin{thmpart}[the Gibbs path --- Theorem B$'$ / B$^\sharp$]\label{thm:B}
Assume EX-A, EX-S, SR, $\phi=1$. For a particle at a state $x\in\Omega\setminus\{\Sz\}$ at the start of $H$ terminal rounds at $\phi=1$
without resampling in between, with the exact data functions $c_*(x)$, $v(x)$ and $\eta(x)$ of the Supplement ($v(x)$ is a two-phase potential,
$\le|M(x)|+a$ while a truth is missing, and $|P|$ on supersets):
\[
\Pr_x(\Sz\text{ not scored within }H\text{ rounds})\ \le\ \Pr\big(\Bin(H,c_*(x))<v(x)\big)+H\eta(x)\ \le\ v(1-c_*)^{\lfloor H/v\rfloor}+H\eta,
\]
with $c_2\ge1/(B+3)-\varepsilon_{\rm del}$ deterministically on supersets, and $c_1\ge1/|M(x)|$-type bounds from pointwise addability (no factorial).
Terminal-round trajectories of different particles are conditionally independent given the terminal-start population and past (IND-T, under EX-S), so the
population failure probability is the product of the per-particle bounds for particles selected using the past only (Cor.\ B-M). At Hamming distance one
the truth is scored in the next round with probability at least $a_1(x)$ (Cor.\ P1).
\end{thmpart}
\noindent\emph{Proof:} see Supplement \S S.5 (source branches \texttt{theory/W1-T24G} \S4, \texttt{theory/W2-VT24G} Lemma F$'$, B1$'$/B2$'$, \texttt{theory/W2-PREM} IND-T).
\emph{Status caveat.} The exact definitions of $c_*(x)$, $v(x)$, $\eta(x)$ are those of W1-T24G \S4 as repaired; W6-FIXREV1 (which corrects Proposition U's
classification or charging) had not delivered its package when this draft was finalised; the statement of $v$ above is therefore the source's inequality form, and the exact constants must be re-read
from the W6-FIXREV1 package before submission.

% STATUS: PROVED-UNDER / PROOF SKETCH / FALSE (as listed); VERIFIED-BY: none beyond W2-VT24G on the NL0(a) row; W2-PREM single run unreviewed; A-REV1: Lemma U finite expansion VERIFIED (limiting constant 0.52 is an ERROR, 0.51), IND-T VERIFIED, "W1 §6 premises imply Hη→0" GAP; SOURCE: W2-PREM:REPORT.md (final); F-REFINE:CARDS.md:Card 4 part D (asymptotic premises)
\begin{remark}[Asymptotic premises of Theorem~\ref{thm:B}]\label{rem:Bpremises}
The asymptotic status of the inputs to Theorem~\ref{thm:B} is mixed and must not be summarised as ``$H\eta\to0$''.
$\tau_2$, $\mathrm{NL0}(0)$, $\mathrm{UDEL}(1)$: PROVED. $\mathrm{NL0}(a)$/$\mathrm{UDEL}(a\ge2)$: PROVED \emph{under} a partner-graph condition PG and FALSE
without it (equicorrelated nulls). TS: PROVED under non-cancellation, FALSE under a cancellation pair. IND-M: PROVED under subset-start concentration and a
non-degeneracy condition (SCR + GS), OPEN for superset starts. On LH all premises are PROOF SKETCH (two population inequalities only numerical). The older
uniform exit hypothesis H24 is \emph{retired}: at boundary states of $\cR_d$ a look-alike of a missing truth exits with probability $\ge(1-o(1))/(B+d+4)$ (W1-T24G CE1,
verified as a deterministic bound; the broad asymptotic wording for correlated designs is a GAP per A-REV1, to be replaced by a narrow statement), and the form
$(1-c')^{\lfloor H/(s_0+d)\rfloor}$ of the brief is invalid (CE2).
\end{remark}

\subsection{Theorem 4: HPM composition (T21)}

% STATUS: PROVED (set inclusions only; "Lemma RO"); VERIFIED-BY: none (single run; Lemma RO is trivial inclusions; W5-LOCKBOOK referee not completed as of f3c037d); SOURCE: W2-BOOK:STATEMENTS.tex:Lemma (Run/output split)
\begin{lemma}[Run/output split]\label{lem:RO}
Under GUARD, for an output-available event $\cO=\cG^c\cap\bar{\cO}$ with $\bar{\cO}\in\sigma(\bar R)$, a reported output $g$ and its potential version $\bar g$:
$\Ind_{\cO}g=\Ind_{\cO}\bar g$; $\Pr(\cO^c\cup\{g\notin A\})\le\Pr(\cG)+\Pr(\bar{\cO}^c)+\Pr(\bar{\cO},\bar g\notin A)$; and
$\E[\Ind_{\cO}g]\le\E[\Ind_{\bar{\cO}}\bar g]$ for $g\ge0$. Identities in expectation are not preserved under $\Pr(\cdot\mid\cO)$.
\end{lemma}
\noindent\emph{Proof:} see Supplement \S S.6 (source branch \texttt{theory/W2-BOOK}, \texttt{PROOF.md} \S0.2).

% STATUS: PROVED (composition: strict maximiser + Markov, locked T21 (C1) restated); route (a) PROVED-LOCKED (T11(iii)); route (d) PROVED; route (e′): drop-in VERIFIED as an implication, route inputs conditional; VERIFIED-BY: W2-VT24G §6 (the three (e′) edits), A-REV1 (Cor B-M VERIFIED under IND-T, "does not prove the presence of eligible particles"); SOURCE: W2-BOOK:STATEMENTS.tex:Theorem (Composition for the reported HPM) [label thm:composition]; W2-BOOK:PROOF.md:§4; F-REFINE:CARDS.md:Card 5
\begin{theorem}[Composition for the reported HPM, T21]\label{thm:T21}
Assume EX-A, EX-S, CAL, SR and GUARD. Let $U\supseteq\cG$ be the no-output event, $\bar U$ its potential-run part, $\bar{\cS}$ the potential scored set and
$\delta_1=1-\pi(\Sz)$. Then, for \emph{every} $n$ and every data set,
\[
\Pr(U\text{ or }\Sh\ne\Sz)\ \le\ \Pr(U)+2\,\E_0\delta_1+\Pr(\bar U^c,\ \Sz\notin\bar{\cS}),\qquad \Pr(U)\le\Pr(\cG)+\Pr(\bar U),
\]
equivalently $\Pr(\cO\cap\{\Sh\ne\Sz\})\le2\E_0\delta_1+\Pr(\bar{\cO},\ \Sz\notin\bar{\cS})$ (Lemma~\ref{lem:RO}). The last term, $\Pr(\bar U^c,\Sz\notin\bar{\cS})=\Pr(\bar{\cO},\Sz\notin\bar{\cS})$, is bounded by each of:
\begin{itemize}
\item[(a)] (i.i.d.\ final draws) $\Pr(\cA^c,\Sz\notin\cS_{\rm cal})+\E\big[\Ind_{\cA}\Ind\{\Sz\notin\cS_{\rm cal}\}(1-\pi_0(\Sz))^{Q_f}\big]$, with $Q_f=4N_e(T+L+5)$;
\item[(d)] (tables and scans) the probability that $\Sz$ lies in no block table and no local search of $\bar R$ is centred at a flip or swap neighbour of $\Sz$;
\item[(d$^+$)] (with Theorem~\ref{thm:T37}) if a data event $\mathrm{TD}^*$ is such that, on $\mathrm{TD}^*$, every completed local search of $\bar R$ centred in $\cR_d$ with $\ell\ge|c\triangle\Sz|$ scores $\Sz$,
then $\Pr(\bar{\cO},\Sz\notin\bar{\cS})\le\Pr_0(\mathrm{TD}^{*c})+\Pr(E_{\rm LS}^c\cap\mathrm{TD}^*)+\Ind\{\ell<s_0+d\}$, where $E_{\rm LS}$ is the event that some local search of $\bar R$ has its centre in $\cR_d$ (drop the last term under the extra-scan rule). With LS0 the start-side event $\{\tilde S_0\in\cR_d\}\subseteq E_{\rm LS}$;
\item[(e$'$)] (Gibbs-path repair at a past-measurable selected state $x$) $\Pr(\cA^c,\Sz\notin\cS_{\rm cal})+\Pr(\cA,\text{no eligible particle})+\E[\Ind_{A_F}\min\{1,\beta_H(X_{\rm sel})\}]$, with
$\beta_H(x)=\Pr(\Bin(H,c_*(x))<v(x))+H\eta(x)$ from Theorem~\ref{thm:B} and $H\le L+5$; and, for eligible particles $i$ at the boundary, $\prod_i\min\{1,\beta_H(X_i)\}$ in place of $\min\{1,\beta_H(X_{\rm sel})\}$ (IND-T).
\end{itemize}
The earlier route (e) with H24 is retired. The premise $\pi(\Sz)>\tfrac12$ is not needed when the local search certifies $\Sz$ (Theorem~\ref{thm:T37}(A3)).
\end{theorem}
\noindent\emph{Proof:} see Supplement \S S.6 (source branches \texttt{theory/W2-BOOK}, \texttt{PROOF.md} \S4 and \texttt{STATEMENTS.tex}; \texttt{theory/W2-VT24G} \S6; locked T21).

% STATUS: PROVED (same inequality, strict-maximiser step re-derived); VERIFIED-BY: same-family referee inside W4-T38FULL; A-REV1 verified the composition inequality of the neighbouring W1-T38E T38(b)-R ("The central composition inequality and liminf manipulation are valid"); SOURCE: W4-T38FULL:T38_STATEMENT.md:Theorem 5 (0)
\begin{remark}[With LS0]\label{rem:T21LS0}
The same inequality holds with $\Spl$, $\cS^+$ and $\bar{\cS}^+$ in place of $\Sh$, $\cS$, $\bar{\cS}$ (the strict-maximiser argument is unchanged):
$\Pr(\cO\cap\{\Spl\ne\Sz\})\le2\E_0\delta_1+\Pr(\bar{\cO}\cap\{\Sz\notin\bar{\cS}^+\})$ (Theorem~\ref{thm:T38}(0)).
\end{remark}

% STATUS: empirical statement (data), not a theorem; VERIFIED-BY: W2-BOOK (counts); SOURCE: W2-BOOK:REPORT.md; F-REFINE:CARDS.md:Card 5 (Data evidence)
\begin{remark}[Reporting $\Pr(U)$ and the no-conditioning rule]\label{rem:PU}
$\Pr(U)$ must be reported empirically: $6/2{,}400=0.25\%$ (Wilson 95\% interval $[0.11\%,0.54\%]$); all six stops are scenario S5; S5 at $p=1000$ shows
2.5\%. The stops are counted as misses. Never condition on completion: a $0.22\%$ stop rate can make $\E[\Zh\mid\text{output}]=0.974\,Z$ (W2-BOOK exact counterexample).
\end{remark}

\subsection{Theorem 5: consistency of the reported HPM (T38) --- the main theorem}

\noindent\textbf{Hypotheses, in groups.}
\emph{Operational:} EX-A, EX-S, CAL, SR, GUARD (Section~\ref{sec:assumptions}); LS0 (Definition~\ref{def:ls0}) is part of the method.
\emph{Posterior:} \textbf{(P)} $\delta_1=1-\pi(\Sz)\to0$ in $\Pr_0$-probability, on a triangular array R1 (sufficient: the hypotheses of Theorem~\ref{thm:T03}).
\emph{Favourable regime:} \textbf{(F)} $1-\pi_0(\Sz)\to0$ in probability (sufficient: (F1), (F2) or (F3) of Theorem~\ref{thm:favourable} at reference $\Sz$).
\emph{Masked regime, route (b1):} there is a \emph{deterministic} reference $S^\circ\in\cR_d$ with $\Phi(S^\circ)$ bounded such that
\begin{itemize}
\item[(M1)] \textbf{ELIG}$(S^\circ)$: $\Pr_0(\pi_0(S^\circ)>\tfrac12)\to1$ (sufficient: G, R-lite, SE-blk, ZL$(S^\circ)$, SPREAD, VIS-lin$(S^\circ)$, by Theorem~\ref{thm:favourable}(F3) at $S^\circ$);
\item[(M2)] $\mathrm{POP\text{-}LAT}(S^\circ,m_n)$: for every $S\in\Lat(S^\circ)$ with $\Phi(S)\ge2$ some progress neighbour $A_S$ of $S$ satisfies $\Delta(B)-\Delta(A_S)\ge m_n$ for every non-progress neighbour $B$ of $S$ and for $B=S$, with $m_n\ge M_n\sqrt{\log n/n}$, $M_n\to\infty$;
\item[(M3)] $\lambda_{\min}(R_{S^\circ\cup\Sz})\ge\varphi_0^2>0$ and $\|\mu\|^2/n\le M$;
\item[(OP-LS0)] LS0 is executed, with its centre emitted, whenever output is available.
\end{itemize}

% STATUS: PROVED for (0), (a), (b1), (b1-LH) (candidates); (b1-LH-iid) PROVED by one unverified run; (c) conditional on T32(i) = CONJECTURE; VERIFIED-BY: (0): W3-VT38E (for W1-T38E's form), W4-T38FULL same-family; A-REV1: "T38(b)-R route (b1), direct ELIG+TD′+OP: VERIFIED"; composition and liminf: valid; Lemma M / LS-score / LH-LAT / Theorem b1: same-family referee only (not reviewed by A-REV1/A-REV2); (b1-LH): inputs VERIFIED (W3-VEXP), assembly unreviewed; (b1-LH-iid): none; SOURCE: W4-T38FULL:T38_STATEMENT.md:Theorem 5 (T38) [to be superseded by W6-FIXREV1:T38_STATEMENT.md if pushed]; F-REFINE:CARDS.md:Card 6
\begin{theorem}[Two-regime consistency of the reported HPM, T38]\label{thm:T38}
Assume EX-A, EX-S, CAL, SR, GUARD and (P) on a triangular array R1 with the v3 budgets indexed by $n$.
\begin{enumerate}
\item[(0)] \textbf{Composition, every $n$.} $\Pr(\cO\cap\{\Spl\ne\Sz\})\le2\E_0[\delta_1]+\Pr(\bar{\cO}\cap\{\Sz\notin\bar{\cS}^+\})$.
Hence, whenever the discovery term tends to $0$, $\Pr(\cO\cap\{\Spl\ne\Sz\})\to0$ and therefore
$\liminf_n\Pr(\cO\cap\{\Spl=\Sz\})\ge1-\limsup_n\Pr(U)$.
\item[(a)] \textbf{Favourable regime.} Under (F) the discovery term tends to $0$, with or without LS0 (the final run's $Q_f\ge4N_e(L+6)$ i.i.d.\ $\pi_0$ draws score $\Sz$; T11(iii)).
\item[(b1)] \textbf{Masked regime} ($\pi_0(\Sz)\to0$ allowed, at the rate of Theorem~\ref{thm:masked}), \textbf{route (b1), LS0.} Under (M1), (M2), (M3), (OP-LS0) and G, on the corresponding events
$\tilde S_0=S^\circ$ and LS0 scores $\Sz$ after at most $\Phi(S^\circ)-1$ moves; the cap $\ell_n\ge\Phi(S^\circ)-1$ suffices. Hence the discovery term tends to $0$.
\item[(b1-LH)] \textbf{An explicit family on which every premise of (b1) holds jointly.} On LH in its exact-Gram realisation (Theorem~\ref{thm:LH}: $p\le n-1$, $|\psi_i|<\min(\rho_{a_i},\rho_{b_i})$, $\psi_i\ne0$, $g\le B_n<2g$, fixed nonzero coefficients), with $S^\circ=\Sz\setminus M$:
(P), (M1), (M2) with a fixed margin $m_0>0$ (Proposition~\ref{prop:LHLAT}) and (M3) all hold; $\pi_0(\Sz)=p^{-c}n^{-c/2}e^{O_P(1)}$; and $Q_f\pi_0(\Sz)\to0$ when $c\ge2$ and $p\to\infty$, and also for $c=1$ when $p/(\sqrt n\log^2n)\to\infty$, so that route (a) fails. With LS0, $\Pr(\cO\cap\{\Spl\ne\Sz\})\to0$ and no statistical input is left open.
\item[(c)] \textbf{Cost.} The number of score evaluations is bounded deterministically by T32(i)'s ceiling plus $(\ell+1)\max_{k\le K}[p+k(p-k)]$ for LS0; under the v3 budgets this is $O((n+p)(p+B)\log^2n+p^2\log^2n)$, polynomial in $n$ under R1. \emph{Conditional on T32(i), which is a CONJECTURE.}
\end{enumerate}
\textbf{Not claimed:} evidence or PIP accuracy in the masked regime; anything given completion, or for fits selected on diagnostics; comparative superiority; any finite-$n$ error rate at $n=200$; a masked $p\gg n$ instance until the single-run LH-iid argument is verified (Theorem~\ref{thm:LH}, iid realisation, is a \emph{candidate input}, not part of this theorem); masked-regime consistency \emph{without} LS0 on any explicit family (OPEN).
\end{theorem}
\noindent\emph{Proof:} see Supplement \S S.7 (source branch \texttt{theory/W4-T38FULL}, \texttt{T38\_PROOF.md} \S\S1--3, 6--7; inputs: Theorems~\ref{thm:T21}, \ref{thm:favourable}, \ref{thm:T37}, Lemma~\ref{lem:LSscore}, Proposition~\ref{prop:LHLAT}).

% STATUS: PROVED-UNDER the cited inputs (candidate); sketch on LH with c=1; VERIFIED-BY: W3-VT38E (Theorem D2 with the repaired constant r*, the finite pilot-stage bound); A-REV1: "Theorem D2 with r* = (11+2√10)/9: VERIFIED as a finite uninterrupted-pilot bound"; "D2-∞: VERIFIED under its additional uniform hypotheses"; "T38(b)-R, entire supplied statement: GAP" (routes (b1′) omit asymptotic rate margins; route (b2) needs an eligible set and the binomial bound; operational qualification must cover pilot and terminal execution); "route (b2), given uniform REPAIR and actual execution: VERIFIED as an implication"; SOURCE: W4-T38FULL:T38_STATEMENT.md:(b1′),(b2); W1-T38E:PROOF.md:§4 D2; W4-SHARE:PROOF.md (envelope hypotheses); F-REFINE:CARDS.md:Card 6
\begin{remark}[Conditional routes --- not part of Theorem~\ref{thm:T38}]\label{rem:routes}
Two further routes would remove the discovery term with \emph{no} method change, but both are conditional and A-REV1 classifies the supplied formulation as a GAP
(to be repaired by W6-FIXREV1). \textbf{(b1$'$) Pilot first stage (D2).} If a one-flip neighbour $x^\circ=\Sz\oplus j$ has $\pi_0(x^\circ)\ge b>1/10$, the dependence parameter satisfies $\mathrm{Dep}\le1-c_1$ off a data-measurable set $\mathcal D$ with $|\mathcal D|\le k_D$, and the no-blow-up event $\mathrm{NB}'$ (effective sample size of the initial pilot weights $\ge0.9H$, or $\Sz$ drawn) has probability $\to1$ (with the corrected constant $r^*=(11+2\sqrt{10})/9$), then the first stage of the tier-0 pilot scores $\Sz$: $\Pr(\Sz\notin\cS_{\rm cal})\le e^{-2Ht^2}+\Pr(\mathrm{NB}'^c)+\exp\{-H(1-\mathrm{Dep})_+g(b-t)/(k_D+B+3)\}$. A statistical instance exists only as a candidate on LH with $c=1$ (D2-a,b,c PROVED by W3-PXN, unverified; NB: PROOF SKETCH); $c\ge2$ is OPEN. \textbf{(b2) Gibbs path.} With a data-measurable eligible set and checkable envelope hypotheses (potentials with $(K_th_t)\ge(1-\eta_t)h_{t-1}$, confinement sets with $\Pr(E^c)\le\varepsilon_E$, fitness oscillation $\Gamma+\sum_t\eta_t\le\log(N\pi_0(h_0)/T_e)-\omega(1)$) and (REP) $\sup_{x\in W}\beta_{L+5}(x)\to0$, then $\Pr(\Sz\notin\cS\mid\cC)\le[\varepsilon_E+\exp\{-N\pi_0(h_0)(1-e^{-c})\}]^4+\varepsilon_{\rm rep}$ (W4-SHARE, PROVED, unverified; no explicit family has these inputs verified: OPEN). The earlier ``bridge mass $\Rightarrow$ delivery'' route is false (C1, COUNTEREXAMPLE, verified by W3-VT38E with a reversible kernel).
\end{remark}

\subsection{Theorem 6: evidence and PIP consistency in the favourable regime (T26(i)--(ii))}

% STATUS: PROVED (composition of locked T14(i), T15, T23(ii-a,b), T22(iv-b), T09(ii), T12(i.a,i.b)); VERIFIED-BY: none final (W2-T26 unverified; W5-LOCKBOOK, branch f3c037d, has referee notes in progress with no final report; status not upgraded); SOURCE: W2-T26:PROOF.md:Theorem A (T26(i)–(ii), premise form), Lemma 1; F-REFINE:CARDS.md:Card 7
\begin{theorem}[Evidence and PIP consistency, T26(i)--(ii)]\label{thm:T26}
Assume EX-A, EX-S, CAL, START, OUT (Remark~\ref{rem:outstart}), and
\begin{itemize}
\item[$(\Omega)$] $\Pr_0(s_0\le K_n)\to1$;
\item[(P1)] for every $\varepsilon>0$, $\Pr\big(\cG_n\cap\{(T_n+1)\Delta_n>\varepsilon\}\big)\to0$ (joint probability over data and calibration streams; $T_n$ the \emph{realised} stage count; sufficient deterministic-budget form $(T^{\max}_n+1)\Delta_n\to0$ in $\Pr_0$-probability), with $\Delta_n=\max\{1-\pi_0(\Sz),1-\pi(\Sz)\}$ (Theorem~\ref{thm:T22});
\item[(P1-ex)] $\Pr_0(\cE_n\cap\{\delta_1>\varepsilon\})\to0$ for every $\varepsilon>0$ (the exact-branch twin of (P1)).
\end{itemize}
Here $\cG_n$ is the event that the SMC branch is taken and the calibration succeeds, $\cE_n$ the exact-enumeration branch. No statistical assumption is used by this theorem. Then, \emph{never conditioning on completion}:
\begin{enumerate}
\item[(A-0)] (fixed $n$, calibration-conditional) for every island $r$: $\E_{\cC}|\Zh^\circ_r/Z-1|\le b(\cC):=2(T+1)\Delta_n/(1-\Delta_n)$; the occupancy fraction off $\Sz$ at the end of stage $T$ and at the start of each measured round has $\E_{\cC}f\le\min\{1,2\delta_1+2b(\cC)\}$; $\E_{\cC}\max_j|\PIPh^\circ_{rj}-\Ind\{j\in\Sz\}|\le\min\{1,2\delta_1+2b(\cC)\}+\bar\rho+u_\pi$ with $\bar\rho\le(B+3)u_\pi$ and $u_\pi=\delta_1/\pi(\Sz)$; $\Pr_{\cC}(\widehat{\MPM}^\circ\ne\Sz)\le2[\min\{1,2\delta_1+2b(\cC)\}+\bar\rho+u_\pi]$;
\item[(A-i)] for every $\varepsilon>0$, $\Pr(\cO_n\cap\{|\Zh/Z-1|>\varepsilon\})\to0$ (also islandwise and for $\log\Zh$);
\item[(A-ii)] for every $\varepsilon>0$, $\Pr(\cO_n\cap\{\max_j|\PIPh_j-\Ind\{j\in\Sz\}|>\varepsilon\})\to0$ and $\Pr(\cO_n\cap\{\widehat{\MPM}\ne\Sz\})\to0$.
\end{enumerate}
The frozen second premise $(L_n+4)\rho^*_n\to0$ is redundant. \textbf{Not claimed:} $\Pr(\cO_n)\to1$; any statement in the masked regime; the secondary ($\Zh_r$-weighted) PIP; the HPM (Theorem~\ref{thm:T38} replaces the old T26(iii)); uniform Monte Carlo rates in $N$; anything about runs selected by their flags or completion; floating-point validity.
\end{theorem}
\noindent\emph{Proof:} see Supplement \S S.8 (source branch \texttt{theory/W2-T26}, \texttt{PROOF.md} \S\S1--3).

% STATUS: PROVED as a composition, CONDITIONAL on T03 + (F1) (hence on the removal of EXP′ by Theorem 1 of W1-EXP); Prop. C (ORTH family) PROOF SKETCH; VERIFIED-BY: none; SOURCE: W2-T26:PROOF.md:Theorem B, Prop. C; F-REFINE:CARDS.md:Card 7 (Theorem B discharged)
\begin{remark}[Where the premise comes from, and its cost]\label{rem:T26B}
In R1 under the hypotheses of Theorem~\ref{thm:T03} and (F1) of Theorem~\ref{thm:favourable}, (P1) holds in its deterministic-budget form, $(T^{\max}_n+1)\Delta_n=O_P\big(\log^2n\,[2^B\sqrt{\log n}\,n^{-\kappa}+e_n+\sqrt{\log n/n}]\big)\to0$ (with $e_n=o_P(n^{-m})$ for all $m$); the ORTH family ($X^\top X=(n-1)I$, $p\le n-1$, fixed $\theta$) satisfies all premises jointly (PROOF SKETCH). Premise (P1) must be stated with the \emph{realised} $T_n$ (random), not with $T^{\max}=4{,}000$: with the deterministic cap the quantity is vacuous below $n\sim10^8$--$10^9$.
\end{remark}

\subsection{Proposition 2: frozen-run identities and the reported bound (T12, T31(i))}

% STATUS: PROVED-LOCKED (T12, T31(i): wave 1, full chain incl. consultant review) + W2-BOOK restatement PROVED (new content: Lemma RO, ¾-covariance formula, output-level bound); VERIFIED-BY: none beyond the locked chain (W5-LOCKBOOK referee not completed as of f3c037d); exact-law propagation reproduces the locked numbers to 1e-50 (W2-BOOK); SOURCE: W2-BOOK:STATEMENTS.tex:Proposition (Frozen-run identities and the reported bound) [label prop:frozen]; W2-BOOK:PROOF.md:§1; F-REFINE:CARDS.md:Card 8
\begin{proposition}[Frozen-run identities and the reported bound]\label{prop:frozen}
Assume EX-A, EX-S, START and CAL, with the calibration $\cC$ (tier, grid, $T$, $N_e$, $L$, round and island counts, local-search cap, inherited score leaders and every
other setting read by the final run) a function of the data and calibration streams only, and escalation decided before any final randomness. For the potential run, on the
success event $\cA$ and for every island $r$ and variable $j$:
\begin{enumerate}
\item[(a)] $\E[\Zh_r\mid\cC]=Z$ and $\E[\Zh_r\PIPh_{rj}\mid\cC]=Z\,\PIP_j$;
\item[(b)] $\E[\Zh\mid\cC]=Z$ and $\E[\tfrac14\sum_r\Zh_r\PIPh_{rj}\mid\cC]=Z\,\PIP_j$;
\item[(c)] $\E[\Zh\,\bar P_j\mid\cC]=Z\,\PIP_j-\tfrac34\mathrm{Cov}(\PIPh_{1j},\Zh_1\mid\cC)$ and $\E[\bar P_j\mid\cC]-\PIP_j=-\mathrm{Cov}(\PIPh_{1j},\Zh_1\mid\cC)/Z$, where $\bar P_j=\tfrac14\sum_r\PIPh_{rj}$ is the reported PIP; both are non-zero in general (so ``pooled evidence $\times$ pooled PIP'' is \emph{not} an identity);
\item[(d)] $\Pr(Z\ge\alpha\Zh\mid\cC)\ge1-\alpha$ for $0<\alpha<1$ (also islandwise), and $\E[\log\Zh\mid\cC]\le\log Z$.
\end{enumerate}
Under GUARD, for the reported evidence, $\Pr(\cO_E,\ \alpha\Zh>Z)\le\alpha\Pr(\cA)$ and $\E[\Ind_{\cO_E}\Zh]\le Z\Pr(\cA)$. If moreover the floating-point evidence satisfies
$|\log\Zh^{\rm fp}-\log\Zh|\le\eta$ outside an event of probability $\varepsilon$ within $\cO_E$ (the validation contract VAL-FP$(\eta,\varepsilon)$), then
$\Pr(\cO_E,\ \alpha\Zh^{\rm fp}>Z)\le\alpha e^{\eta}+\varepsilon$. The secondary ($\Zh_r$-weighted) PIP is removed.
\end{proposition}
\noindent\emph{Proof:} see Supplement \S S.9 (source branch \texttt{theory/W2-BOOK}, \texttt{PROOF.md} \S1; locked T12, T31(i)).

% STATUS: COUNTEREXAMPLE (exact, mpmath 50 digits); empirical; VERIFIED-BY: W2-BOOK (self); VAL-FP: not measured (needs an R/C++ replay); SOURCE: W2-BOOK:PROOF.md:§0.3; W2-BOOK:REPORT.md; F-REFINE:CARDS.md:Card 8
\begin{remark}[What is not claimed]\label{rem:frozenno}
No unbiasedness given completion is claimed: $\E[\Zh^{\rm rep}\mid\cO_E]=Z$ is false in general (W2-BOOK exact counterexample: $0.22\%$ stops give $\E[\Zh\mid\text{output}]=0.974\,Z$). The reported PIP is not unbiased for $\PIP_j$ (part (c)). The VAL-FP contract has not been measured.
\end{remark}

% =====================================================================================================
\section{What the theory does and does not say at $n=200$}\label{sec:disclosure}
% =====================================================================================================

% STATUS: disclosure; every numerical fact below is quoted with its source and status; VERIFIED-BY: as listed per item; SOURCE: F-REFINE:ARCHITECTURE.md:§6 (overstatement watch), §7 (n = 200 disclosure); F-REFINE:CARDS.md:§9, Cards 1,3,4,6; W2-VT03R:REPORT.md:R1,R2
\textbf{The theorems are statements about mechanism and about $n\to\infty$; they are not statements about the reported numbers at the study's size
($n=200$, $p\in\{500,1000\}$, $s_0=14$).} The simulations are evidence in their own right. The following facts must be read with every theorem above.

\subsection{What the theory does say}

% STATUS: PROVED exact inequalities (non-asymptotic); VERIFIED-BY: per theorem (see source comments above); SOURCE: F-REFINE:ARCHITECTURE.md:§6 item 4
The only non-vacuous \emph{finite-$n$} statements are exact inequalities: the composition (C1) of Theorem~\ref{thm:T21} (including its form with LS0, Remark~\ref{rem:T21LS0}), the
deterministic Theorem~\ref{thm:T37} (with the lattice form of (A3), and Lemma~\ref{lem:LSscore}), the ESS-share inequality of the pilot analysis, Corollary P1 of Theorem~\ref{thm:B} at $B=4$, and Proposition~\ref{prop:frozen}.
In particular, the decomposition~\eqref{eq:story} is exact, and at the study's size the error is carried by its first term (item 2 of the list in Section~\ref{sec:disclosure}).

\subsection{What the theory does not say}

\begin{enumerate}
\item[1.] \textbf{The start at $n=200$ is dilution-masked, not cancellation-masked.}
% STATUS: PROVED (exact, W3-VEXP Prop. D; W3-PXN Prop. W); thresholds PROOF SKETCH; VERIFIED-BY: single runs; SOURCE: F-REFINE:CARDS.md:Card 3 (Data evidence); W3-VEXP:REPORT.md:Prop. D; W3-PXN:PROOF.md:Prop. W
The block tables score against the full response, so each truth sees the other truths as noise ($v_T\approx7$--$26\sigma^2$); the start mode misses $6$--$12$ of the $14$ truths at every signal scale, and this disappears only for $n$ of order $800$--$1600$ in these designs. The masked start of Theorems~\ref{thm:masked} and~\ref{thm:LH} is a different mechanism (cancellation across blocks), and LH's cancellation cannot survive the study's partition (groups of 5 with $B=14$: co-blocking kills LH masking, W3-PXN Prop.~Q, PROOF SKETCH; LH needs $g\le B<2g$). The paper must therefore present Theorems~\ref{thm:masked}--\ref{thm:LH} and~\ref{thm:T38}(b1) as the masked regime the sampler is designed for, \emph{not} as the study's masking.

\item[2.] \textbf{Every completed wrong HPM is a posterior miss.}
% STATUS: empirical (THEORY_PRIORITIES §5); VERIFIED-BY: n/a; SOURCE: F-REFINE:CARDS.md:Card 1 (Data evidence), Card 5; F-REFINE:ARCHITECTURE.md:§7(ii)
In the $p\gg n$ study, $565$ of $2{,}394$ completed fits have $F(\Sz)\le F(\Sh)$ ($0$ discovery misses); there are $6$ resource stops, all in S5, counted as misses ($6/2{,}400=0.25\%$). So the binding term at this size is the posterior term $2\E_0\delta_1$ of~\eqref{eq:story}; the repair statements (Theorems~\ref{thm:T37}, \ref{thm:W}) must be read ``whenever the posterior allows''.

\item[3.] \textbf{Theorem~\ref{thm:T03} is asymptotic, with proof thresholds far above $n=200$.}
% STATUS: PROVED (threshold computation; the earlier 1.2e5 claim was an ERROR); VERIFIED-BY: W2-VT03R (R1, numbers reproduced), A-REV2 (§3.5 confirms; W6-FIXREV2 independent recomputation d213e7b confirms per its commit message, final report pending); SOURCE: W2-VT03R:REPORT.md:R1, R2
The deterministic steps of the proof, as written, hold only from $n\approx2.1\times10^6$ ($p=n^{1.17}$) to $2.5\times10^6$ ($p=n^{1.30}$): the binding condition is $q_n\le\tfrac12$ (the earlier figure $1.2\times10^5$ is an ERROR, corrected by W2-VT03R and confirmed by A-REV2); the $O_P$ constants are unspecified. The sparse-eigenvalue condition $\mathrm{SE}_w$ ranges over at most $114$ columns at $n=200$ \emph{because the cap $K=100$ binds}, not because of the exact size tail (uncapped $m^\sharp=146/168$); it is unverified on the realised designs (greedy-adversarial eigenvalue probes: $\lambda_{\min}(G)\le0.40$ for S1, $\le0.22$ for S4 at $|\Sz\cup J|=114$; feasibility of any fixed $g_0$ is unverified). By interlacing, it requires every $100$ columns of $X$ to be linearly independent. The beta-min constant of the proof exceeds the executed signal by at least $43\times$ (proof needs $\theta^{*2}\ge27.5/\varphi_0^2$ at $p=500$; executed $\theta^{*2}=0.64/0.16$), and the orthogonal-design scale $(\log n+2\log p)/n\approx0.09$ is $7\times$ \emph{below} the executed signal, so the shortfall is the constant, not the signal.

\item[4.] \textbf{The start rates are vacuous at the study size.}
% STATUS: PROVED (explicit arithmetic on the rate formulas); VERIFIED-BY: none (arithmetic on F-REFINE's quoted formulas; not independently recomputed here); SOURCE: F-REFINE:CARDS.md:Card 3 (Data evidence, "Rate relevance"); NOTE: F-REFINE:ARCHITECTURE.md:§7(iv) says the factor is "≥ 10^3 at B = 14 for n ≤ 10^4"; the CARDS arithmetic is used here because it is the stated product
The rates of Theorem~\ref{thm:favourable} carry a factor $2^B$: $2^B\sqrt{\log n/n}$ at $B=14$ equals $2.7\times10^3$ at $n=200$, $497$ at $n=10^4$, $61$ at $n=10^6$, so they are numerically vacuous below $n\approx10^9$. (Only the $2^B$-free iid-Gaussian rate of Theorem~\ref{thm:LH}, a single unverified run, avoids this.) Theorem~\ref{thm:masked} bounds i.i.d.\ draws from $\pi_0$ only, and says nothing about the hybrid sampler.

\item[5.] \textbf{The HPM at this size is found by Gibbs flips, a route with no completed theorem.}
% STATUS: empirical + OPEN; VERIFIED-BY: none (counts from the campaign's regenerated study data; the E_LS extraction is open); SOURCE: F-REFINE:CARDS.md:Card 4 (Data evidence), Card 6 (Data evidence)
$98\%$ of fresh scores in the study come from Gibbs flips; the final local search repaired $0$ fits; whether the study's $541$ final centres lay in $E_{\rm LS}$ is unextracted (OPEN). The paper proves exact finite-$n$ inequalities for this route (Theorem~\ref{thm:B}) but no rate: the single-particle bound is vacuous at $n=200$ ($H\eta\approx4$--$16$; $v>H$ possible). The route (b2) of Remark~\ref{rem:routes}, which would match the observed behaviour, has no verified instance. The local search is not the delivery mechanism at $n=200$, and ``the local search repairs'' must not appear as an explanation of the numbers.

\item[6.] \textbf{LS0 does not change the picture at $n=200$.}
% STATUS: empirical (post hoc, W3-PXN) + OPEN (O-LS0 data check pending); VERIFIED-BY: none (O-LS0 pushed scripts only); SOURCE: F-REFINE:CARDS.md:Card 6 (Data evidence); F-REFINE:ARCHITECTURE.md:§4; main:theory_campaign/HQ/SR_AUDIT.md
LS0 cannot fix posterior misses; the start mode is $8$--$14$ moves from $\Sz$ at $n=200$ (not a bounded $\Phi(S^\circ)$), the cap $\ell=10$--$11$ binds in S4/S6, and post hoc on regenerated study data LS0 scored $\Sz$ in $113/128$ fits with $|\tilde S_0\triangle\Sz|\le\ell+1$ and $13/32$ beyond. The approval of LS0 is conditional on the O-LS0 data check (few or no changed reported HPMs, never turning a correct report into an incorrect one beyond posterior-mode corrections, negligible compute), whose result was not available at drafting. LS0 buys the one masked-regime consistency statement whose premises are all proved on an explicit family (Theorem~\ref{thm:T38}(b1-LH)); it buys nothing for route (b2) or for evidence/PIP claims in the masked regime.

\item[7.] \textbf{Evidence and PIP consistency are favourable-regime, long-hypothesis statements.}
% STATUS: PROVED (as composition) with the relevance caveat; VERIFIED-BY: none (see Theorem 6); SOURCE: F-REFINE:CARDS.md:Card 7 (Data evidence, Risk); F-REFINE:ARCHITECTURE.md:§6 item 9
At $n=200$, $\Delta_n\approx1$ (dilution) so Theorem~\ref{thm:T26} is silent, although the sampler still estimates $Z$ within $1$--$2\%$ in the masked design. (P1) must be stated with the realised $T_n$.

\item[8.] \textbf{Statements the paper must not make.}
% STATUS: constraint list (not a theorem); VERIFIED-BY: n/a; SOURCE: F-REFINE:ARCHITECTURE.md:§6 items 1--14 (items not already covered above are listed)
(i) The literal T37/$\mathrm{TD}(d,m)$, ``exactly $|S\triangle\Sz|$ moves'', the $2^{s_0}$ union bound, ``T37 replaces GAIN, LK, EXIT'' (true only for the local-search route and only on $E_{\rm LS}$); (ii) H24 or any uniform-over-$\cR_d$ exit hypothesis, the $(1-c')^{\lfloor H/(s_0+d)\rfloor}$ form, T24's factorial; (iii) ``the exact BB tail makes SE meaningful at $n=200$'' (the cap does, and only for $\mathrm{SE}_w$; anchoring is cosmetic, Remark~\ref{rem:T03}); (iv) any finite-$n$ error rate at $n=200$ from Theorems~1, 2, 3(ii), 5 or 6; (v) ``the correlated-design advantage follows from T37's margin'' (contradicted at $n=200$ by the $c=0.5$ posterior misses in S4--S6; the LH margin of Proposition~\ref{prop:LHLAT} is the mechanism on LH, with no statement about the study); (vi) ``T27 is a lower bound for the method'' (it bounds i.i.d.\ start draws only); (vii) ``T26 with budgets indexed by $n$'' via $T^{\max}=4{,}000$; (viii) $\E[\Zh\mid\text{completed}]=Z$ and ``pooled $\Zh\times$ pooled PIP is unbiased'' (both false); (ix) the cost corollary as a theorem (T32(i) is a CONJECTURE); (x) the word ``VERIFIED'' as if it meant reviewed: all verification by the campaign is same-family (Claude checking Claude) except the two cross-family reviews A-REV1 and A-REV2, which cover only the items whose verdicts are quoted above; (xi) ``D2 covers one-cancellation masking at $p\gg\sqrt n$'' (conditional; its CF instances are not masked); (xii) ``no method change needed'': true for everything except Theorem~\ref{thm:T38}(b1), which is the headline and rests on LS0.

\item[9.] \textbf{Implementation facts the theorems lean on have had at most one reading.}
% STATUS: IMPL; VERIFIED-BY: HQ read of step_gibbs for SR only (SR_AUDIT.md); none for GUARD, IND-T, scan-before-limit, LS0 emit_centre; SOURCE: F-REFINE:CARDS.md:§9 item 5; main:theory_campaign/HQ/SR_AUDIT.md; W5-LOCKBOOK:AUDIT.md (f3c037d, unreviewed)
SR (HQ's read of \texttt{step\_gibbs}: current and flipped configurations are evaluated before the accept draw; no call site disables emission), GUARD (stop-only; code reading of \texttt{v3\_runtime.h}, \texttt{v3\_engine.h}, \texttt{v3\_core.cpp}), the scan-before-limit rule of the local search, IND-T's sequential stream reading, and LS0's \texttt{emit\_centre}. None has had an independent audit confirmed on these branches.
\end{enumerate}

\end{document}
